\section{The Bethe--Salpeter equation from the ladder}\label{app:D.12}

In the bubble of \cref{sec:4.7,sec:4.8} the electron and the hole propagate independently between the two density vertices. The ladder diagrams let them scatter off each other any number of times in between, and the Bethe--Salpeter equation of \cref{sec:4.9} is their sum. This section derives it with the rules of \cref{app:D.11}. The derivation follows the many-body treatments of Hanke and Sham~\cite{hanke1980many} and of Strinati~\cite{strinati1988application}; the first-principles form is that of Refs.~\cite{albrecht1998ab,rohlfing2000electron,onida2002electronic}.

\subsection*{Transitions and the pair propagator}

Take a basis of orbitals $\psi_m(\vect{r})$ that diagonalize the one-body part of the problem, with energies $\xi_m$ measured from $\mu_{\mathrm{ch}}$ and the spin in the label. In a crystal $m$ stands for a band, a crystal momentum and a spin. The density operator is $n_e(\vect{q}) = \sum_{mn}\rho_{mn}(\vect{q})\,c^\dagger_mc_n$ with $\rho_{mn}(\vect{q}) = \langle m|\ee^{\ii\vect{q}\cdot\vect{r}}|n\rangle$, by \cref{eq:2.3-onebody}. Each term moves one electron from $n$ to $m$, and we call the pair $(mn)$ a transition. The density response is then a combination of the correlation functions of the pair operators $c^\dagger_mc_n$. Define, in imaginary time,
\begin{equation}
  \mathfrak{L}_{mn,m'n'}(\tau) = -\bigl\langle\mathrm{T}_\tau\,c^\dagger_n(\tau)c_m(\tau)\;c^\dagger_{m'}(0)c_{n'}(0)\bigr\rangle ,
  \label{eq:D.12-L}
\end{equation}
the pair propagator, written in fraktur; it is the two-particle correlation function with both creation and both annihilation times set equal in pairs. Since $n_e(\vect{q})^\dagger = \sum_{mn}\rho_{mn}(\vect{q})^*c^\dagger_nc_m$, the response of \cref{eq:4.7-Rtau} with $O_a = n_e(\vect{q}+\vect{G})$ and $O_b = n_e(\vect{q}+\vect{G}')$ is
\begin{equation}
  R_{\vect{G}\vect{G}'}(\vect{q},\ii\omega^{\mathrm{B}}_j) = \frac1\Vol\sum_{mn,m'n'}\rho_{mn}(\vect{q}+\vect{G})^*\,\mathfrak{L}_{mn,m'n'}(\ii\omega^{\mathrm{B}}_j)\,\rho_{m'n'}(\vect{q}+\vect{G}') ,
  \label{eq:D.12-R}
\end{equation}
where we take $\vect{q}$ not equal to a reciprocal-lattice vector, so that $\langle n_e(\vect{q}+\vect{G})\rangle = 0$ and only connected contractions contribute. Everything the target does to the pair is in $\mathfrak{L}$.

For independent quasiparticles, Wick's theorem pairs $c_m(\tau)$ with $c^\dagger_{m'}(0)$ and $c_{n'}(0)$ with $c^\dagger_n(\tau)$, which forces $m' = m$ and $n' = n$. The two contractions form one closed loop through the two pair operators, and by the counting of \cref{app:D.11} the result is $\GF_m(\tau)\GF_n(-\tau)$. Its transform is the frequency sum \cref{eq:D.10-sum},
\begin{equation}
  \mathfrak{L}^0_{mn,m'n'}(\ii\omega^{\mathrm{B}}_j) = \delta_{mm'}\delta_{nn'}\,T\sum_{j'}\GF_m(\ii\omega^{\mathrm{F}}_{j'} + \ii\omega^{\mathrm{B}}_j)\,\GF_n(\ii\omega^{\mathrm{F}}_{j'})
  = \delta_{mm'}\delta_{nn'}\,\frac{n_{\mathrm{F}}(\xi_n) - n_{\mathrm{F}}(\xi_m)}{\ii\omega^{\mathrm{B}}_j - (\xi_m - \xi_n)} ,
  \label{eq:D.12-L0}
\end{equation}
diagonal in transitions, with a pole at the transition energy $\xi_m - \xi_n$ and the Pauli factor $n_{\mathrm{F}}(\xi_n) - n_{\mathrm{F}}(\xi_m)$. With plane waves in \cref{eq:D.12-R} it is the bubble \cref{eq:4.7-bubble}.

\subsection*{First-order contractions}

The interaction is the Coulomb operator of \cref{sec:2.3} in the orbital basis. Under $\mathrm{T}_\tau$, with the creation operators at $\tau + 0^+$, it is a product of two bilinears, as in \cref{app:D.11},
\begin{equation*}
  H_{\mathrm{int}} = \frac12\sum_{abcd}v_{ab,cd}\,\bigl[c^\dagger_ac_c\bigr]\bigl[c^\dagger_bc_d\bigr], \qquad
  v_{ab,cd} = \int\dd^3r\,\dd^3r'\,\psi_a^*(\vect{r})\psi_c(\vect{r})\,v_{\mathrm{C}}(\vect{r},\vect{r}')\,\psi_b^*(\vect{r}')\psi_d(\vect{r}') ,
\end{equation*}
with $v_{\mathrm{C}}(\vect{r},\vect{r}') = e^2/(4\pi|\vect{r} - \vect{r}'|)$ the transform of \cref{eq:A.2-vC}; spin is conserved at each bilinear. The rules of \cref{app:D.11} were derived for plane waves, but their proof used only this bilinear structure, and they hold in any orbital basis with $v_{ab,cd}$ in place of $v_{\mathrm{C}}(k)/\Vol$. The expansion \cref{eq:D.11-Gexp}, applied to \cref{eq:D.12-L}, gives at first order $+\int_0^\beta\dd\tau_1\,\langle\mathrm{T}_\tau\,c^\dagger_nc_m(\tau)\,H_{\mathrm{int}}(\tau_1)\,c^\dagger_{m'}c_{n'}(0)\rangle_0$, connected part only. The two bilinears of $H_{\mathrm{int}}$ can be contracted with the external ones in three topologically different ways, which we call the exchange, the direct and the self-energy contractions.

In the exchange contraction, each bilinear of $H_{\mathrm{int}}$ closes a loop with one external pair operator: $[c^\dagger_ac_c]$ with $c^\dagger_nc_m(\tau)$, which needs $a = m$, $c = n$, and $[c^\dagger_bc_d]$ with $c^\dagger_{m'}c_{n'}(0)$, which needs $b = n'$, $d = m'$. Bilinears are even in fermion operators and can be regrouped at no cost, and the average factorizes into two free pair averages, each equal to $-\mathfrak{L}^0$. The product is $+\mathfrak{L}^0\mathfrak{L}^0$. The two choices of which bilinear goes with which end cancel the $\tfrac12$, and the term is $+\mathfrak{L}^0_{mn}(\tau - \tau_1)\,v_{mn',nm'}\,\mathfrak{L}^0_{m'n'}(\tau_1)$, a convolution in $\tau$. In the direct contraction, both bilinears lie on the single loop that runs through the two pair operators. One of them is on the electron line, which enters at $c^\dagger_{m'}$ and leaves at $c_m$, and requires $a = m$, $c = m'$. The other is on the hole line and requires $b = n'$, $d = n$. The diagram has the same four propagators and the same interaction line as the exchange contraction, but one closed loop instead of two, and by the loop rule of \cref{app:D.11} it carries the opposite sign: $-\mathfrak{L}^0_{mn}(\tau - \tau_1)\,v_{mn',m'n}\,\mathfrak{L}^0_{m'n'}(\tau_1)$. In the self-energy contractions, both bilinears of $H_{\mathrm{int}}$ act on one propagator, or one of them closes a loop by itself; these are the Hartree and Fock self-energies of \cref{app:D.11} inserted into the electron or the hole line, and they shift the orbital energies.

The exchange and direct contractions define the kernels of the Bethe--Salpeter equation:
\begin{equation}
  \begin{aligned}
  K^{\mathrm{x}}_{\mathrm{full},mn,m'n'} &= \int\dd^3r\,\dd^3r'\;\psi_m^*(\vect{r})\psi_n(\vect{r})\;v_{\mathrm{C}}(\vect{r},\vect{r}')\;\psi_{m'}(\vect{r}')\psi_{n'}^*(\vect{r}') ,\\
  K^{\mathrm{d}}_{mn,m'n'} &= \int\dd^3r\,\dd^3r'\;\psi_m^*(\vect{r})\psi_{m'}(\vect{r})\;W^{\mathrm{scr}}(\vect{r},\vect{r}')\;\psi_{n'}^*(\vect{r}')\psi_n(\vect{r}') .
  \end{aligned}
  \label{eq:D.12-kernels}
\end{equation}
The exchange kernel $K^{\mathrm{x}}_{\mathrm{full}}$ carries the whole Coulomb interaction, including the long-range term that the transition-basis section below splits off. It is the term of the random-phase approximation, and its name refers to the exchange of the electron and the hole between the two transitions. In the direct kernel $K^{\mathrm{d}}$ the electron stays at $\vect{r}$ and scatters from $m'$ to $m$, while the hole stays at $\vect{r}'$ and scatters from $n'$ to $n$. It is written above already with the screened interaction that the next step supplies. A diagonal element $K^{\mathrm{d}}_{mn,mn} = \int\!\int|\psi_m(\vect{r})|^2W^{\mathrm{scr}}(\vect{r},\vect{r}')|\psi_n(\vect{r}')|^2$ is positive, and it enters with a minus sign.

\subsection*{The statically screened ladder}

The Bethe--Salpeter equation keeps all diagrams built by repeating these two insertions, with the following choices, which define the approximation. The orbital energies $\xi_m$ are the quasiparticle energies of \cref{sec:4.7}, so that the self-energy contractions are summed into the propagators, with the quasiparticle widths neglected and $z_{\mathrm{qp}}$ set to one. A direct interaction line is screened: a chain of bubbles inserted into it turns $v_{\mathrm{C}}$ into $v_{\mathrm{C}} + v_{\mathrm{C}}Rv_{\mathrm{C}} = \varepsilon^{-1}v_{\mathrm{C}} = W^{\mathrm{scr}}$, the screened interaction of \cref{sec:4.8}. An exchange line is not screened, since a bubble chain on it is already counted as repeated exchange. These choices are consistent with the $GW$ self-energy of \cref{sec:4.8}: the kernel is the change of the Hartree potential plus the $GW$ self-energy when $\GF$ changes, with the change of $W^{\mathrm{scr}}$ itself neglected~\cite{hanke1980many,strinati1988application}. The Hartree term gives $K^{\mathrm{x}}_{\mathrm{full}}$ and the $GW$ term gives $-K^{\mathrm{d}}$.

A screened direct line carries a frequency. In the direct diagram the electron line has the fermionic frequency $\ii\omega^{\mathrm{F}}_{j_1} + \ii\omega^{\mathrm{B}}_j$ on one side of the line and $\ii\omega^{\mathrm{F}}_{j_2} + \ii\omega^{\mathrm{B}}_j$ on the other, the hole line correspondingly, and the line is $W^{\mathrm{scr}}$ at the bosonic difference $\ii\omega^{\mathrm{F}}_{j_1} - \ii\omega^{\mathrm{F}}_{j_2}$. The two internal frequency sums are then coupled through the screened line, and $\mathfrak{L}$ does not close on itself: the pair splits into an electron and a hole at different times. The static approximation replaces $W^{\mathrm{scr}}(\ii\omega^{\mathrm{B}})$ by $W^{\mathrm{scr}}(\omega = 0)$ at all frequencies, an interaction that is instantaneous in $\tau$. With it, the diagram with one direct line is
\begin{equation*}
  \begin{aligned}
  &T\sum_{j_1}\GF_m(\ii\omega^{\mathrm{F}}_{j_1} + \ii\omega^{\mathrm{B}}_j)\GF_n(\ii\omega^{\mathrm{F}}_{j_1})\;\bigl(-K^{\mathrm{d}}_{mn,m'n'}\bigr)\;T\sum_{j_2}\GF_{m'}(\ii\omega^{\mathrm{F}}_{j_2} + \ii\omega^{\mathrm{B}}_j)\GF_{n'}(\ii\omega^{\mathrm{F}}_{j_2}) \\
  &\qquad = \mathfrak{L}^0_{mn}\,\bigl(-K^{\mathrm{d}}_{mn,m'n'}\bigr)\,\mathfrak{L}^0_{m'n'} ,
  \end{aligned}
\end{equation*}
since each sum separately is \cref{eq:D.10-sum}. The exchange line is instantaneous and carries $\ii\omega^{\mathrm{B}}_j$. Every further direct line repeats the step, and the sum of the ladder is the geometric series of \cref{eq:C.7-dyson},
\begin{equation}
  \mathfrak{L}(\ii\omega^{\mathrm{B}}_j) = \mathfrak{L}^0(\ii\omega^{\mathrm{B}}_j) + \mathfrak{L}^0(\ii\omega^{\mathrm{B}}_j)\,\bigl[K^{\mathrm{x}}_{\mathrm{full}} - K^{\mathrm{d}}\bigr]\,\mathfrak{L}(\ii\omega^{\mathrm{B}}_j) ,
  \label{eq:D.12-BSE}
\end{equation}
a matrix equation in the space of transitions at a single frequency. This is the Bethe--Salpeter equation in the static ladder approximation, whose range of validity is stated in \cref{sec:4.9}.

\subsection*{Spin structure of the kernel}

For a spin-unpolarized target the orbitals come in pairs $m\!\uparrow$, $m\!\downarrow$ with the same spatial part. The exchange kernel couples a transition of either spin to a transition of either spin, because the density at $\vect{r}'$ sums both. The direct kernel conserves the spin of the electron and of the hole separately, and with $\psi_m^*\psi_n$ requiring equal spins it couples only transitions of the same spin. In the combination $(mn\!\uparrow + mn\!\downarrow)/\sqrt2$ the kernel is therefore $g_{\mathrm{deg}}K^{\mathrm{x}}_{\mathrm{full}} - K^{\mathrm{d}}$, with $g_{\mathrm{deg}} = 2$ and the kernels now written for spatial orbitals; in $(mn\!\uparrow - mn\!\downarrow)/\sqrt2$ it is $-K^{\mathrm{d}}$. These are the singlet and triplet excitons. The density vertex is the same for both spins, so $n_e(\vect{q})$ creates only singlets, with the weight $\sqrt{g_{\mathrm{deg}}}$ per transition. From here on the labels are spatial, the kernel is the singlet one, and $g_{\mathrm{deg}}$ appears in front of the response.

\subsection*{The transition-basis Hamiltonian}

At $T = 0$, in an insulator, the Pauli factor of \cref{eq:D.12-L0} is $+1$ for a transition from an occupied band $v$ to an empty band $c$. These are the resonant transitions; the reverse ones, with the factor $-1$, are antiresonant, and we mark them with a bar. Transitions between two occupied or two empty orbitals have $\mathfrak{L}^0 = 0$ and drop out. In the matrix notation of \cref{app:C.7}, let $\mathbf{F}$ be the diagonal matrix of the Pauli factors, with $\mathbf{F}^2 = \mathbf{1}$, and $\mathbf{D}$ the diagonal matrix of the transition energies; the inverse of \cref{eq:D.12-L0} is then $\mathbf{F}(\ii\omega^{\mathrm{B}}_j - \mathbf{D})$. \Cref{eq:D.12-BSE} reads $\mathfrak{L}^{-1} = (\mathfrak{L}^0)^{-1} - \mathbf{K}$ with $\mathbf{K} = g_{\mathrm{deg}}K^{\mathrm{x}}_{\mathrm{full}} - K^{\mathrm{d}}$, and
\begin{equation}
  \mathfrak{L} = \bigl[\mathbf{F}(\ii\omega^{\mathrm{B}}_j - \mathbf{D}) - \mathbf{K}\bigr]^{-1} = \bigl(\ii\omega^{\mathrm{B}}_j - \mathbf{D} - \mathbf{F}\mathbf{K}\bigr)^{-1}\mathbf{F} .
  \label{eq:D.12-Lsol}
\end{equation}
The poles are the eigenvalues of $\mathbf{D} + \mathbf{F}\mathbf{K}$, which in blocks of resonant and antiresonant transitions is
\begin{equation}
  \mathbf{D} + \mathbf{F}\mathbf{K} = \begin{pmatrix} H^{\mathrm{exc}}_{\mathrm{full}} & K_{\mathrm{r}\bar{\mathrm{r}}} \\ -K_{\mathrm{r}\bar{\mathrm{r}}}^* & -H^{\mathrm{exc}*}_{\mathrm{full}}\end{pmatrix}, \qquad
  H^{\mathrm{exc}}_{\mathrm{full},cv,c'v'} = (\xi_c - \xi_v)\,\delta_{cc'}\delta_{vv'} + g_{\mathrm{deg}}K^{\mathrm{x}}_{\mathrm{full},cv,c'v'} - K^{\mathrm{d}}_{cv,c'v'} .
  \label{eq:D.12-blocks}
\end{equation}
The lower blocks follow from \cref{eq:D.12-kernels}. Exchanging $m$ and $n$ in both transitions conjugates $K^{\mathrm{x}}_{\mathrm{full}}$, since $v_{\mathrm{C}}$ is real. It also conjugates $K^{\mathrm{d}}$ when $W^{\mathrm{scr}}(\vect{r},\vect{r}') = W^{\mathrm{scr}}(\vect{r}',\vect{r})$ is real, as the static screened interaction of a target with time-reversal symmetry is. The factor $-1$ of $\mathbf{F}$ on the antiresonant rows supplies the signs. The same two properties make $H^{\mathrm{exc}}_{\mathrm{full}}$ Hermitian. In the resonant block the exchange kernel carries the singlet factor $g_{\mathrm{deg}}$~\cite[Eqs.~(18) and (20)]{blase2020bethe}.

The coupling block $K_{\mathrm{r}\bar{\mathrm{r}}}$ connects the creation of a pair with its annihilation. In a semiconductor its effect on the eigenvalues is of second order in the kernel divided by twice the gap, and the Tamm--Dancoff approximation drops it.

In a crystal the Coulomb line of $K^{\mathrm{x}}_{\mathrm{full}}$ carries $\vect{q} + \vect{G}$ for every reciprocal-lattice vector $\vect{G}$, and its term at $\vect{G} = 0$, which we call $K^{\mathrm{x}}_0$, is the long-range macroscopic field. The rest is the exchange kernel proper, $K^{\mathrm{x}} = K^{\mathrm{x}}_{\mathrm{full}} - K^{\mathrm{x}}_0$, the local-field kernel. We write $H^{\mathrm{exc}}$ for $H^{\mathrm{exc}}_{\mathrm{full}}$ with $K^{\mathrm{x}}$ in place of $K^{\mathrm{x}}_{\mathrm{full}}$. With the band and crystal-momentum labels of a crystal written out, a transition at fixed $\vect{q}$ takes an electron from $v\vect{k}$ to $(c,\vect{k}+\vect{q})$, and the resonant problem is \cref{eq:4.9-bse}, a Hermitian eigenvalue problem, whose eigenvalues $E_\lambda$ and envelopes $A^\lambda$ depend on $\vect{q}$. The same problem with $H^{\mathrm{exc}}_{\mathrm{full}}$ has the eigenpairs $E^{\mathrm{full}}_\lambda$ and $A^\lambda_{\mathrm{full}}$. In the crystal the kernels \cref{eq:D.12-kernels} become the sums over reciprocal-lattice vectors of \cref{eq:4.9-kernels}, with the convention of \cref{app:A.3}, and the full exchange kernel is
\begin{equation}
  K^{\mathrm{x}}_{\mathrm{full},cv\vect{k},c'v'\vect{k}'} = K^{\mathrm{x}}_{cv\vect{k},c'v'\vect{k}'} + \frac1\Vol\,v_{\mathrm{C}}(q)\,\rho_{cv}(\vect{k},\vect{q})\,\rho_{c'v'}(\vect{k}',\vect{q})^* .
  \label{eq:D.12-crystal}
\end{equation}
The direct kernel follows from \cref{eq:D.12-kernels} with the expansion
\begin{equation*}
  W^{\mathrm{scr}}(\vect{r},\vect{r}') = \Vol^{-1}\sum_{\vect{p}\vect{G}\vect{G}'}\ee^{\ii(\vect{p}+\vect{G})\cdot\vect{r}}\,W^{\mathrm{scr}}_{\vect{G}\vect{G}'}(\vect{p})\,\ee^{-\ii(\vect{p}+\vect{G}')\cdot\vect{r}'} :
\end{equation*}
the integral over $\vect{r}$ gives $\langle c,\vect{k}+\vect{q}|\ee^{\ii(\vect{p}+\vect{G})\cdot\vect{r}}|c',\vect{k}'+\vect{q}\rangle$, nonzero for $\vect{p} = \vect{k} - \vect{k}'$, and the integral over $\vect{r}'$ the complex conjugate of $\langle v\vect{k}|\ee^{\ii(\vect{p}+\vect{G}')\cdot\vect{r}}|v'\vect{k}'\rangle$. No symmetry of the crystal is needed. The exchange kernels carry the momentum $\vect{q}$ of the pair and the direct kernel the momentum $\vect{p}$ exchanged between the electron and the hole. All three scale as $1/\Vol = 1/(N\Omega)$, and the sum over $\vect{k}'$ in \cref{eq:4.9-bse} makes $H^{\mathrm{exc}}$ act as an integral operator of order one, as the count of \cref{app:A.3} requires.

\subsection*{Response, structure factor and dielectric function}

Continued to $\omega + \ii0^+$ (\cref{eq:4.7-Rcont}), the resonant block of \cref{eq:D.12-Lsol} in the Tamm--Dancoff approximation is the resolvent of the pair Hamiltonian, $(\omega - H^{\mathrm{exc}} + \ii0^+)^{-1} = \sum_\lambda A^\lambda A^{\lambda\dagger}/(\omega - E_\lambda + \ii0^+)$, and likewise for $H^{\mathrm{exc}}_{\mathrm{full}}$. For the problem \cref{eq:4.9-bse}, \cref{eq:D.12-R} with the singlet weight gives
\begin{equation}
  \bar R_{00}(\vect{q},\omega) = \frac{g_{\mathrm{deg}}}{\Vol}\sum_\lambda\frac{\bigl|\sum_{cv\vect{k}}A^\lambda_{cv\vect{k}}\,\rho_{cv}(\vect{k},\vect{q})^*\bigr|^2}{\omega - E_\lambda(\vect{q}) + \ii0^+} + (\text{antiresonant}) ,
  \label{eq:D.12-Rexc}
\end{equation}
where the antiresonant part has its poles at $-E_\lambda$. Since $K^{\mathrm{x}}$ lacks the term at $\vect{G} = 0$, this is the head of the response that is proper with respect to the long-range Coulomb term. The complete response follows from it, because the difference of the two kernels, $g_{\mathrm{deg}}K^{\mathrm{x}}_0$ with $K^{\mathrm{x}}_0 = \Vol^{-1}v_{\mathrm{C}}(q)\rho(\vect{q})\rho(\vect{q})^\dagger$ in transition space, has rank one. Let $\bar{\mathfrak{L}}$ solve \cref{eq:D.12-BSE} with $K^{\mathrm{x}}$ in place of $K^{\mathrm{x}}_{\mathrm{full}}$. Then $\mathfrak{L}^{-1} = \bar{\mathfrak{L}}^{-1} - g_{\mathrm{deg}}K^{\mathrm{x}}_0$, which is $\mathfrak{L} = \bar{\mathfrak{L}} + \bar{\mathfrak{L}}\,g_{\mathrm{deg}}K^{\mathrm{x}}_0\,\mathfrak{L}$, and its head in \cref{eq:D.12-R} closes on itself:
\begin{equation*}
  R_{00} = \bar R_{00} + \bar R_{00}\,v_{\mathrm{C}}(q)\,R_{00}, \qquad R_{00} = \frac{\bar R_{00}}{1 - v_{\mathrm{C}}(q)\bar R_{00}} .
\end{equation*}
The same $R_{00}$ is \cref{eq:D.12-Rexc} with the eigenpairs $E^{\mathrm{full}}_\lambda$ and $A^\lambda_{\mathrm{full}}$ of the full kernel. It is the complete density response, the one whose imaginary part enters the dark-matter rate. At $T = 0$ and $\omega > 0$ only the resonant part contributes to $\Imag R_{00}$, and the fluctuation--dissipation theorem \cref{eq:4.2-fdt} gives the exciton contribution to the structure factor,
\begin{equation}
  S(\vect{q},\omega) = \frac{2\pi g_{\mathrm{deg}}}{\Vol}\sum_\lambda\Bigl|\sum_{cv\vect{k}}A^\lambda_{\mathrm{full},cv\vect{k}}(\vect{q})\,\rho_{cv}(\vect{k},\vect{q})^*\Bigr|^2\,\delta\bigl(\omega - E^{\mathrm{full}}_\lambda(\vect{q})\bigr) .
  \label{eq:D.12-S}
\end{equation}

The macroscopic dielectric function needs the proper head. Since $(\varepsilon^{-1})_{00} = 1 + v_{\mathrm{C}}R_{00}$ (\cref{sec:4.4}) and $1 + v_{\mathrm{C}}R_{00} = 1/(1 - v_{\mathrm{C}}\bar R_{00})$,
\begin{equation}
  \varepsilon_{\mathrm{mac}}(\vect{q},\omega) = \frac{1}{(\varepsilon^{-1})_{00}} = 1 - v_{\mathrm{C}}(q)\,\bar R_{00}(\vect{q},\omega) .
  \label{eq:D.12-epsmac}
\end{equation}
The block inversion \cref{eq:C.7-head} does the same for the dielectric matrix of the random-phase approximation; here the local fields of \cref{sec:4.6} and the electron--hole attraction are treated in one equation. In the optical limit, with $\rho_{cv} \to \vect{q}\cdot\vect{v}_{cv}/(\varepsilon_{c\vect{k}} - \varepsilon_{v\vect{k}})$ by \cref{eq:3.2-optical}, the imaginary part is the $\varepsilon_2$ of \cref{eq:4.9-eps2}.
